\honest{An Alien God}{Eliezer Yudkowsky}{2007}

``A curious aspect of the theory of evolution,'' said Jacques Monod, ``is that everybody thinks he understands it.''

Nature looks full of purpose: a rabbit's legs are built for running, a fox's jaws for tearing. Before Darwin, I grant, God was a reasonable explanation, and Paley inferred a watchmaker from a watch. But look at all of nature. Foxes are well designed to catch rabbits, and rabbits to escape foxes; was the Creator unable to make up Its mind? I would not design a toaster with one part sending electricity to the coils and another blocking it. Even the cactus, which might seem made to give desert animals water, is covered in spines. A horde of competing gods, Hindu or Shinto, would explain this better than one. Did anyone ever notice the evidence for Hinduism over Christianity? ``Probably not.'' \nb{In Hume's \textsc{Dialogues Concerning Natural Religion} (1779), Philo asks why ``several deities'' might not have framed the world, as many workers build a ship.} Much of nature is also cruel. Darwin could not believe that a beneficent God had designed ichneumon wasps to feed ``within the living bodies of Caterpillars,'' and I wonder whether any earlier thinker saw in this the evidence for Manichaean religions. \nb{Hume's Philo did, eighty years before Darwin's letter: the Manichaean system ``has more probability than the common hypothesis,'' though he then prefers first causes with ``neither goodness nor malice.'' That is close to the conclusion of this post.}

The answer, we all know, is ``evolution.'' But I worry that people absorb evolution as a magical purpose factory, as in \textsc{X-Men}, where Storm gets the power to throw lightning in one mutation, and stronger radiation means stronger mutations. A theory is strong by what it prohibits, and evolution could not explain a toaster. \nb{The principle is Karl Popper's; see ``Your Strength as a Rationalist.''}

George Williams observed that many non-biologists think rattles grow on rattlesnakes for their own benefit. That kind of purpose is not allowed. Evolution runs on a correlation between how genes build organisms and how many copies of those genes reach the next generation; there is no ``Evolution Fairy'' deciding which genes would be helpful. People who accept this still ask which genes are ``helpful,'' as if a rattlesnake's gene could help non-rattlesnakes. The rattle spread because of the rattle, probably because snakes with better rattles survived more often: perhaps predators avoid stepping on them, or the rattle draws a dog's bite to the tail instead of the head. A gene becomes more common when its effects cause more copies of it to reach the next generation. A gene that costs one unit and gives three to a sibling, who shares half one's genes, may spread.

This explains the conflict: there are as many evolutions as populations. It explains the cruelty: anesthetizing the wasp's prey, or letting an old elephant die painlessly, would cost evolution nothing, but ``There's no one to argue with.'' Humans make up justifications for what they want; there is no Fairy of elephant evolution that wants what is best for elephants and argues for it before an overseer who cares about fitness. There is no advocate for the elephants anywhere in the system. Genes that replicate more become more common, like water flowing downhill, and equally benevolent.

Humans who imagine better designs and then argue that they would raise fitness predict wrongly. An engineer would make taste track our needs, so that lettuce tasted delicious to the obese; but calories were reliably scarce for our ancestors, so evolution made us love them. We are the embodied history of which organisms did in fact survive and reproduce, not of which ought to have. The retina is wired backward, with the nerves in front of the light-sensitive cells and a blind spot where they exit, while other animals evolved it the right way round. No single mutation can turn the whole retina around; turning the cells without rewiring the nerves leaves the animal blind. An engineer can change many parts at once and plan ahead. Evolution cannot, and is as blind as a half-redesigned retina.

Theologians were nearer the truth than believers in spontaneous generation: evolution is ``closer to God than it is to pure random entropy.'' Mutation is random, but selection is not: a small statistical correlation between a gene and its bearer's reproduction adds up, over millions of years, to something very powerful. Damien Broderick said that gods must be ontologically distinct from creatures, and evolution is not a creature. Like God it is bodiless, everywhere and unmade. But it has no mind, designs badly, is divided against itself, and is not nice. In a way Darwin discovered God, but not the one theology expected. Had he found a bodiless mind that loves us, people would have said ``That's God!'' Evolution is not a god, but if it were, it would not be Jehovah but Lovecraft's Azathoth, the blind idiot god at the center of everything.

So much for religionists who claim to believe in a vague deity; had they meant it, they would have recognized their creator when Darwin found it. \nb{The post names none of them. It also says itself that ``Evolution is not a God'' and that ``the Maker has no mind.'' A believer who declined to call evolution a god was agreeing with the post.} So much, too, for those who say they are waiting, curious, for science to discover God: it has already found the godlike maker of humans, but they wanted their own specific God. I am glad it was Azathoth and not Odin. I like having a Creator I can outwit. It beats being a pet.
